% Provenance (never printed): see sources/outline_sources.md (Year 3:
% partition coefficients; batch and fractional melting; crystallisation
% models; reading a spider diagram quantitatively). Brief:
% sources/book5/BRIEFS.md, ch. 12. Curves from figdata/bachelor-3/
% melting_models.py. Exercise and problem values in the Book 5 scratch
% calculator calc/ch12.py. Aldren Island and its lavas are invented.
\chapter{النمذجة بالعناصر النزرة}\label{ch:b3:trace-element-modelling}

تفجّر بازلتان على الجزيرة نفسها على بُعد بضعة كيلومترات أحدهما من الآخر. يحوي أحدهما
أكثر من مرة ونصف قدر ما يحويه الآخر من اللانثانوم، وما يقارب أربع مرات
أقل من الإتربيوم. فهل انصهرا من وشاحين مختلفين، أم من وشاح واحد
انصهر بمقادير مختلفة، أم عند أعماق مختلفة؟ لا تكشف العناصر الرئيسية
ذلك إلا بالكاد، إذ إنّ كليهما بازلت. أمّا العناصر النزرة، فهي تتبع المعادن
دون أن تغيّرها، فتحفظ سجلًّا لكمية المصهور الذي أُنتج وكيفية خروجه
من مصدره والمعادن التي بقيت خلفه. يكتب هذا الفصل ذلك السجلّ
في صورة معادلات ثم يقرؤه.

\begin{recall}
معامل توزيع المعدن والمصهور $D$، ومعامل التوزيع الكلّي،
والعناصر المتوافقة وغير المتوافقة، والعناصر الأرضية النادرة، والكوندريت CI
والوشاح البدائي، والمخططات المُطبَّعة، وشذوذ اليوروبيوم
وMORB وOIB موجودة في مجلد السنة الثانية، في فصله عن العناصر في
الأرض. يرد البيريدوتيت واللهرزوليت والانصهار الجزئي وكسر المصهور $F$
في فصله عن انصهار الوشاح؛ والتبلّر التوازني والتبلّر التجزيئي،
وقانون رايلي للتجزّؤ، والامتزاج،
والتبلّر التجزيئي المصحوب بالاستيعاب (AFC)، في فصله عن
التمايز الصهاري. وفيه أُثبت ثبات نسبة عنصرين شديدي عدم التوافق
أثناء التبلّر.
\end{recall}

\section{حصيلة صخرة في طور الانصهار}\label{sec:b3:trace-element-modelling:balance}

تنصهر صخرة كتلتها 1. عند كسر المصهور $F$ تكون قد أنتجت كتلة مقدارها $F$ من
المصهور، وأبقت كتلة مقدارها $1 - F$ من البقايا الصلبة. يتوزّع عنصر نزر
موجود بالتركيز $C_0$ في المصدر بين الطورين، ويكون تركيزه
في المصهور $C_L$ وفي البقايا $C_S$. كل نموذج في هذا الفصل صورة من حصيلة واحدة
هي $C_0 = FC_L + (1 - F)C_S$ للمصهور والبقايا معًا،
وقاعدة تربط بين $C_S$ و$C_L$: عند التلامس بين البلّورة والمصهور،
أي $C_{\text{معدن}} = D\,C_L$. تختلف النماذج في ما يلامس ماذا
ومتى.

تتفاوت معاملات توزيع معادن الوشاح تفاوتًا كبيرًا. ففي
دراسة تجميعية واحدة، يبلغ معامل توزيع اللانثانوم $D = 0.058$ في الكلينوبيروكسين و$0.001$ في
الجارنيت والأوليفين والأورثوبيروكسين والسبينل، ويبلغ معامل توزيع الإتربيوم $D = 0.58$ في
الكلينوبيروكسين و$4.5$ في الجارنيت و$0.16$ في المعادن الثلاثة الأخرى. يأخذ الجارنيت
العناصر الأرضية الثقيلة ويرفض الخفيفة.
% ledger: budahn91:d-ree

\begin{definition}[الانصهار غير الشكلي]\label{def:b3:trace-element-modelling:nonmodal}
في \emph{الانصهار غير الشكلي}\index{الانصهار غير الشكلي}، تدخل معادن المصدر إلى
المصهور بنسب $p_i$ تختلف عن كسورها الكتلية
$x_i$ في المصدر. يرى المصهور حينئذ معامل التوزيع الكلّي
للمصدر، $D_0 = \sum_i x_iD_i$، في بداية الانصهار، ثم
المعامل المرجَّح بالانصهار $P = \sum_i p_iD_i$ كلّما تقدّم الانصهار.
\end{definition}

نادرًا ما يكون الانصهار نمطيًا بالكامل. في مثال محلول على لهرزوليت سبينلي
يحوي \qty{58}{\percent} من الأوليفين و\qty{27}{\percent} من الأورثوبيروكسين،
\qty{12}{\percent} من الكلينوبيروكسين و\qty{3}{\percent} من السبينل، وتدخل المعادن
إلى المصهور بالنسب 20 و25 و45 و\qty{10}{\percent}:
ينصهر الكلينوبيروكسين أسرع من حصته بأربع مرات، والأوليفين
أبطأ من حصته بثلاث مرات.
% ledger: white13:ex73

\begin{proposition}[معامل التوزيع الكلّي للبقايا]
\label{prop:b3:trace-element-modelling:residue}
في \omterm{def:b3:trace-element-modelling:nonmodal}{الانصهار غير الشكلي} عند ثبات $p_i$، تملك البقايا عند كسر المصهور $F$
معامل التوزيع الكلّي التالي
\[
\bar D(F) = \frac{D_0 - PF}{1 - F},
\]
صالح ما دام كلّ معدن باقيًا، أي عندما يكون $F < x_i/p_i$ لكل $i$.
\end{proposition}

\begin{proof}
لكل وحدة كتلة من المصدر، تكون كتلة المعدن $i$ المتبقية عند $F$ هي $x_i - p_iF$،
وكتلة البقايا $1 - F$. ومعاملها الكلّي هو المتوسط الموزون بالكتلة
المتوسط $\sum_i (x_i - p_iF)D_i/(1 - F) = (D_0 - PF)/(1 - F)$. يُعدّ المعدن مستنفدًا
عندما يكون $x_i - p_iF = 0$.
\end{proof}

بالاستعانة بمعاملات المعادن أعلاه، فإنّ اللهرزوليت السبينلي من المثال المحلول،
واللهرزوليت الجارنيتي المحفوظ عند \qty{5}{GPa}
و\qty{1400}{\celsius} في المختبر، ويحوي \qty{54.1}{\percent} أوليفين،
\qty{9.8}{\percent} أورثوبيروكسين، و\qty{19.5}{\percent} كلينوبيروكسين
و\qty{16.7}{\percent} جارنيت بالكتلة، له معاملات التوزيع الكلّية الواردة أدناه.
% ledger: white13:ex73, rfr25:gkr-5gpa, budahn91:d-ree
\begin{center}
\small
\begin{tabular}{lccc}
\hline
المصدر & $D_{\mathrm{La}}$ & $D_{\mathrm{Sm}}$ & $D_{\mathrm{Yb}}$ \\
\hline
لهرزوليت سبينلي & 0.0078 & 0.062 & 0.21 \\
لهرزوليت جارنيتي & 0.012 & 0.107 & 0.97 \\
\hline
\end{tabular}
\end{center}
اللانثانوم شديد عدم التوافق في الحالتين. والإتربيوم متوسط عدم التوافق
في اللهرزوليت السبينلي ومحايدًا تقريبًا في اللهرزوليت الجارنيتي. ويكمن
الفرق بين المصدرين كلّه تقريبًا في العناصر الأرضية الثقيلة
النادرة.

\section{الانصهار الدفعي}\label{sec:b3:trace-element-modelling:batch}

\begin{theorem}[الانصهار الدفعي]\label{thm:b3:trace-element-modelling:batch}
في \emph{الانصهار الدفعي}\index{الانصهار الدفعي}، يبقى المصهور في توازن
مع البقايا كلّها حتى يخرج دفعةً واحدة. عند ثبات معامل التوزيع الكلّي
$D$، يكون للمصهور والبقايا عند كسر المصهور $F$
القيمتان التاليتان:
\[
\frac{C_L}{C_0} = \frac{1}{D + F(1 - D)}, \qquad
\frac{C_S}{C_0} = \frac{D}{D + F(1 - D)}.
\]
\end{theorem}

\begin{proof}
عند التوازن تحمل كلّ بلّورة من البقايا $D_iC_L$، ومنه تحمل البقايا
$C_S = DC_L$ في المتوسط. يعطي توازن الكتلة $C_0 = FC_L + (1 - F)DC_L$
قيمة $C_L$، ويُستنتج منها $C_S$.
\end{proof}

المعادلة هي قانون التبلور عند التوازن من مجلد السنة الثانية مقروءًا
بالعكس: ففي الحالتين، $F$ هو كسر السائل في نظام مغلق. ويستحقّ حدّاها
الاطلاع عليهما: عندما $F \to 0$ يكون $C_L \to C_0/D$، أي أنّ أوّل قطرة من المصهور
لا تُغني العنصر إلّا بمقدار أقصاه $1/D$. وعندما $D \to 0$ يكون $C_L \to
C_0/F$: يُخفَّف العنصر غير المتوافق تمامًا في المصهور، ويتحدّد
إغناؤه بالكسر $F$ وحده. أمّا العنصر المتوافق، أي $D > 1$، فيكون المصهور
أفقر من المصدر، وتزداد فقرًا كلّما صَغُر $F$.

\begin{proposition}[الانصهار الدفعي غير الشكلي]\label{prop:b3:trace-element-modelling:nonmodal}
في \omterm{thm:b3:trace-element-modelling:batch}{الانصهار الدفعي} بنسب الانصهار $p_i$ المعرّفة في \cref{def:b3:trace-element-modelling:nonmodal}،
\[
\frac{C_L}{C_0} = \frac{1}{D_0 + F(1 - P)}.
\]
\end{proposition}

\begin{proof}
تحمل البقايا $\bar D(F)\,C_L$ (انظر \cref{prop:b3:trace-element-modelling:residue})،
ومنه $C_0 = FC_L + (1 - F)\bar D(F)C_L = C_L(F + D_0 - PF)$.
\end{proof}

عندما يتجاوز $P$ القيمة $D_0$، كما في حالة اللانثانوم في المثال المحلول، حيث $D_0 =
0.0074$ و$P = 0.025$، تفقد البقايا أفضل حاملٍ لها، وهو الكلينوبيروكسين، بسرعة
من سواه من المعادن. يكون المصهور عندئذٍ أغنى قليلًا مما يتنبّأ به الانصهار الشكلي،
بأقلّ من \qty{2}{\percent} عند $F = 0.1$ في حالة اللانثانوم (التمرين~6)، لكنّ
البقايا تنفد من الكلينوبيروكسين عند $F = 0.12/0.45 \approx 0.27$، وبعد ذلك
تتغيّر المعادن والمعاملات.
% ledger: white13:ex73

\section{الانصهار التجزيئي والانصهار الديناميكي}\label{sec:b3:trace-element-modelling:fractional}

يتجمّع المصهور في البيريدوتيت على حواف الحبيبات، ويستطيع الإفلات ما إن يصبح
متّصلًا. والطرف المقابل للانصهار الدفعي يزيل كلّ قطرة
لحظة تكوّنها.

\begin{theorem}[الانصهار التجزيئي]\label{thm:b3:trace-element-modelling:fractional}
في \emph{الانصهار التجزيئي}\index{الانصهار التجزيئي}، تكون كلّ زيادة من المصهور
في توازن مع البقايا لحظة تكوّنها فقط، ثم تفارقها.
مع ثبات $D$، وبعد استخراج كتلة كلّية $F$ من المصهور، تحمل البقايا
والمصهور المتكوّن في تلك اللحظة القيمتين التاليتين:
\[
\frac{C_S}{C_0} = (1 - F)^{1/D - 1}, \qquad
\frac{C_L}{C_0} = \frac{1}{D}(1 - F)^{1/D - 1}.
\]
\end{theorem}

\begin{proof}
تملك البقايا كتلة $1 - F$ وتركيزًا $C_S$. وعندما تغادر كتلة
$\mathrm dF$ من المصهور بتركيز $C_L = C_S/D$، ينخفض العنصر في البقايا
بمقدار $\mathrm d[(1 - F)C_S] = -(C_S/D)\,\mathrm dF$. بالنشر،
$(1 - F)\,\mathrm dC_S = C_S(1 - 1/D)\,\mathrm dF$، ومنه
$\mathrm dC_S/C_S = (1/D - 1)\,\mathrm d[-\ln(1 - F)]$. وبالتكامل بدءًا من
$C_S = C_0$ عند $F = 0$ نحصل على $\ln(C_S/C_0) = (1/D - 1)\ln(1 - F)$،
و$C_L = C_S/D$.
\end{proof}

بالنسبة إلى العنصر غير المتوافق، يكون $1/D - 1$ كبيرًا وموجبًا: تُستنزف البقايا
في غضون بضعة أضعاف $D$ من الانصهار، ويهبط المصهور اللحظي
من $C_0/D$ إلى ما يكاد يكون لا شيء. وتكون بقايا \omterm{thm:b3:trace-element-modelling:fractional}{الانصهار التجزيئي} أشدّ
استنفادًا بكثير من بقايا \omterm{thm:b3:trace-element-modelling:batch}{الانصهار الدفعي}. غير أنّ المصهورات لا تُقذف واحدةً
تلو الأخرى، بل تتجمّع في قنوات وتكوّن برَكًا، \omterm{def:b3:trace-element-modelling:fractional}{والمصهور التجزيئي المجمَّع}، المعرَّف أدناه،
يُرسم بخطّ متقطّع في الشكل.

\begin{omfigure}
\begin{tikzpicture}
\begin{groupplot}[group style={group size=2 by 1, horizontal sep=1.6cm},
  omrecord, width=0.47\linewidth, height=0.44\linewidth, title style={font=\small},
  xlabel={كسر المصهور $F$}, xmin=0, xmax=0.3, ymode=log, ymin=0.01, ymax=150,
  xtick={0,0.1,0.2,0.3}, xticklabels={0,0.1,0.2,0.3},
  ytick={0.01,0.1,1,10,100}, yticklabels={0.01,0.1,1,10,100},
  legend style={font=\scriptsize, legend cell align=left}]
\nextgroupplot[title={(أ) الدفعي (متّصل)، المجمَّع (متقطّع)},
  ylabel={$C_L/C_0$}, legend style={at={(0.97,0.97)}, anchor=north east}]
% computed in figdata/bachelor-3/melting_models.py
\addplot[omDef, thick] table[x=f, y=b001] {figdata/out/bachelor-3-melting_models-curves.dat};
\addlegendentry{$D = 0.01$}
\addplot[red!70!black, thick] table[x=f, y=b01] {figdata/out/bachelor-3-melting_models-curves.dat};
\addlegendentry{$D = 0.1$}
\addplot[black, thick] table[x=f, y=b1] {figdata/out/bachelor-3-melting_models-curves.dat};
\addlegendentry{$D = 1$}
\addplot[omLithLime!70!black, thick] table[x=f, y=b10] {figdata/out/bachelor-3-melting_models-curves.dat};
\addlegendentry{$D = 10$}
\addplot[omDef, thick, densely dashed] table[x=f, y=a001] {figdata/out/bachelor-3-melting_models-curves.dat};
\addplot[red!70!black, thick, densely dashed] table[x=f, y=a01] {figdata/out/bachelor-3-melting_models-curves.dat};
\addplot[omLithLime!70!black, thick, densely dashed] table[x=f, y=a10] {figdata/out/bachelor-3-melting_models-curves.dat};
\nextgroupplot[title={(ب) التجزيئي، اللحظي}]
\addplot[omDef, thick] table[x=f, y=i001] {figdata/out/bachelor-3-melting_models-curves.dat};
\addplot[red!70!black, thick] table[x=f, y=i01] {figdata/out/bachelor-3-melting_models-curves.dat};
\addplot[black, thick] table[x=f, y=i1] {figdata/out/bachelor-3-melting_models-curves.dat};
\addplot[omLithLime!70!black, thick] table[x=f, y=i10] {figdata/out/bachelor-3-melting_models-curves.dat};
\end{groupplot}
\end{tikzpicture}

\par\smallskip{\small إغناء العنصر الأثري في المصهور بدلالة
كسر المصهور، لأربع قيم لمعامل التوزيع الكلّي. (أ) مصهورات \omterm{thm:b3:trace-element-modelling:batch}{الانصهار الدفعي} (متّصل)
والمصهورات التجزيئية المجمَّعة (متقطّعة؛ عند $D = 1$ يبقى كلاهما عند 1). (ب)
المصهورات اللحظية للانصهار التجزيئي: عنصر ذو $D = 0.01$ يُستنفد
من البقايا بعد بضعة في المئة من الانصهار.}\label{fig:b3:trace-element-modelling:curves}
\end{omfigure}

\begin{definition}[المصهور التجزيئي المجمَّع]\label{def:b3:trace-element-modelling:fractional}
\emph{المصهور التجزيئي المجمَّع}\index{المصهور التجزيئي المجمَّع} هو
مزيج جميع المصهورات اللحظية للانصهار التجزيئي من $0$ إلى
$F$، مجمَّعةً دون تفاعل لاحق مع البقايا.
\end{definition}

\begin{proposition}[تركيب المصهور المجمَّع]
\label{prop:b3:trace-element-modelling:aggregate}
مع $D$ ثابت، يكون للمصهور التجزيئي المجمَّع
\[
\frac{\bar C_L}{C_0} = \frac{1 - (1 - F)^{1/D}}{F}.
\]
\end{proposition}

\begin{proof}
العنصر المستخرَج هو المصدر مطروحًا منه البقايا،
$C_0 - (1 - F)C_S = C_0[1 - (1 - F)^{1/D}]$ حسب \cref{thm:b3:trace-element-modelling:fractional}،
موزَّعًا على كتلة $F$ من المصهور.
\end{proof}

عندما يصبح $(1 - F)^{1/D}$ صغيرًا، يحتوي المصهور المجمَّع على العنصر كلّه،
و$\bar C_L \approx C_0/F$ مهما كانت $D$: بالنسبة إلى العناصر شديدة عدم التوافق،
تتفق المصهورات التجزيئية المجمَّعة مع مصهورات \omterm{thm:b3:trace-element-modelling:batch}{الانصهار الدفعي} ذات $D$ الصغيرة،
وكلاهما يقيس $F$.

الانصهار الحقيقي يقع بين الطرفين النهائيين. تحتفظ البقايا ببعض المصهور في مساماتها،
والمصهور الذي يغادر هو الفائض فوق هذه المسامية.

\begin{definition}[الانصهار الديناميكي]\label{def:b3:trace-element-modelling:dynamic}
في \emph{الانصهار الديناميكي}\index{الانصهار الديناميكي} (ويُسمّى أيضًا الانصهار
المستمر)، تحتفظ البقايا بنسبة كتلية ثابتة $\phi$ من المصهور في
مساماتها، في حالة توازن مع بلوراتها، ويُستخرَج كلّ مصهور يتشكّل
فوق ذلك.
\end{definition}

\begin{proposition}[الانصهار الديناميكي]\label{prop:b3:trace-element-modelling:dynamic}
\omterm{def:b3:trace-element-modelling:dynamic}{الانصهار الديناميكي} ذو المسامية $\phi$ هو انصهار تجزيئي بمعامل التوزيع الفعّال
$D' = (1 - \phi)D + \phi$، حيث $F$ هو الكسر الكتلي
للمصهور المستخرَج:
\[
\frac{C_L}{C_0} = \frac{1}{D'}(1 - F)^{1/D' - 1}.
\]
\end{proposition}

\begin{proof}
البقايا بلورات بكسر $1 - \phi$ ومصهور مسامي بكسر
$\phi$، وكلاهما في توازن مع المصهور عند $C_L$، فيكون
$C_S = [(1 - \phi)D + \phi]C_L = D'C_L$. ثم تنطبق
\cref{thm:b3:trace-element-modelling:fractional} مع $D'$.
\end{proof}

عند $\phi = 0$ يكون النموذج انصهارًا تجزيئيًا؛ وعند $\phi \to 1$ يقترب
من \omterm{thm:b3:trace-element-modelling:batch}{الانصهار الدفعي}. مسامية مقدارها واحد في المئة تحدّ إثراء المصهور الأول
عند $1/\phi = 100$، مهما كانت $D$ صغيرة: إذ يعمل المصهور المسامي نطاقًا عازلًا
للبقايا.

\begin{omfigure}
\begin{tikzpicture}[font=\small, >=Stealth]
\tikzset{rock/.style={draw, fill=omLithLime!25, minimum width=2.3cm, minimum height=2.0cm},
  melt/.style={draw, fill=red!25}}
% (a) batch
\node[rock] (a) at (0,0) {};
\foreach \x/\y in {-0.7/0.5, 0.3/0.6, -0.2/-0.1, 0.7/-0.4, -0.6/-0.6, 0.4/0.0}
  \fill[red!55] (\x,\y) circle (0.09);
\node[align=center, font=\scriptsize] at (0,-1.45) {\foreignlanguage{arabic}{البقايا $1-F$ عند $DC_L$}\\\foreignlanguage{arabic}{المصهور $F$ عند $C_L$}};
\draw[->, thick, red!70!black] (0,1.05) -- (0,1.7) node[above, font=\scriptsize] {دفعة واحدة};
\node[font=\small] at (0,-2.25) {\foreignlanguage{arabic}{(أ) دفعي}};
% (b) fractional
\node[rock] (b) at (4,0) {};
\draw[->, thick, red!70!black] (4,1.05) -- (4,1.7);
\foreach \y/\s in {1.85/0.06, 2.05/0.08, 2.27/0.1}
  \fill[red!55] (4,\y) circle (\s);
\node[align=center, font=\scriptsize] at (4,-1.45) {كلّ قطرة تغادر\\حين تتشكّل};
\node[font=\small] at (4,-2.25) {\foreignlanguage{arabic}{(ب) تجزيئي}};
% (c) dynamic
\node[rock] (c) at (8,0) {};
\foreach \x/\y in {7.4/0.5, 8.4/0.3, 7.8/-0.5, 8.6/-0.6}
  \fill[red!55] (\x,\y) circle (0.06);
\draw[->, thick, red!70!black] (8,1.05) -- (8,1.7);
\foreach \y/\s in {1.85/0.06, 2.05/0.08, 2.27/0.1}
  \fill[red!55] (8,\y) circle (\s);
\node[align=center, font=\scriptsize] at (8,-1.45) {\foreignlanguage{arabic}{المسام تحفظ $\phi$؛}\\\foreignlanguage{arabic}{يغادر الفائض}};
\node[font=\small] at (8,-2.25) {\foreignlanguage{arabic}{(ج) ديناميكي}};
\end{tikzpicture}

\par\smallskip{\small ثلاث طرق لخروج المصهور من مصدر منصهر. (أ)
دفعي: يبقى المصهور مع البقايا كلّها، ثم يغادر. (ب) تجزيئي:
تغادر كلّ زيادة فورًا؛ والزيادات المجمَّعة تكوّن المصهور المجمَّع.
(ج) ديناميكي: تحتفظ البقايا بكسر كتلي صغير $\phi$ من المصهور في
مساماتها.}\label{fig:b3:trace-element-modelling:modes}
\end{omfigure}

\section{الانصهار، ثم التبلّر}\label{sec:b3:trace-element-modelling:crystallisation}

نادرًا ما يثور المصهور الذي يغادر مصدره دون أن يتغيّر. يتبلّر
الأوليفين أولًا، ثم البيروكسين والبلاجيوكلاز في طريقه إلى الأعلى، وقد يستوعب صخر
الجدار، وقد يختلط بدفعات أخرى. تنطبق قوانين مجلد السنة الثانية
بدورها: الحمم مصهور بكسر $F$ من مصدره، ثم سائل
احتفظ بكسر $f$ من كتلته عبر التبلّر التجزيئي، بحيث
يكون، لعنصر له $D$ في البقايا المنصهرة و$D_c$ في
المعادن المتبلّرة،
\[
C_{\text{الحمم}} = C_0\,\frac{1}{D + F(1 - D)}\,f^{\,D_c - 1}.
\]
العامل الثاني هو نفسه لكلّ عنصر ترفضه المعادن المتبلّرة،
$D_c \approx 0$: فهو يُدرّجها جميعًا بمقدار $1/f$. ومن ثمّ تُحدَّد نسب
هذه العناصر عند المصدر، والحمم التي فقدت الأوليفين
تحتفظ بنسبة La/Yb لمصهورها الأوّلي، بينما يرتفع اللانثانوم فيها. ولاستنتاج $F$
من تركيز يلزم معرفة $f$ أيضًا؛ أمّا استنتاج المصدر من نسبة
فلا يتطلّب ذلك. البلاجيوكلاز هو الاستثناء بين العناصر الأرضية النادرة: يأخذ الأوروبيوم،
وإزالته تترك شذوذ الأوروبيوم في مجلد السنة الثانية.

\section{قراءة أنماط العناصر الأرضية النادرة}\label{sec:b3:trace-element-modelling:patterns}

يتشارك مصهور الليرزوليت السبينيلي ومصهور الليرزوليت الجارنيتي العناصر
الأرضية النادرة الخفيفة، التي لا تحتفظ بها أيٌّ من البقايا، ويختلفان في الثقيلة منها،
التي يحتفظ بها الجارنيت. وعلى مخطط مُطبَّع بالكوندريت، تكون مصهورات الجارنيت
شديدة الانحدار ومسطّحة القاع عند الطرف الثقيل؛ أمّا مصهورات السبينيل فأكثر اعتدالًا. ويزداد
الانحدار كلّما انخفضت $F$.

\begin{omfigure}
\begin{tikzpicture}
\begin{axis}[omrecord, width=0.8\linewidth, height=0.5\linewidth, ymode=log,
  xmin=56.5, xmax=71.5, ymin=1, ymax=1000, xtick={57,58,60,62,63,64,65,70,71},
  xticklabels={La,Ce,Nd,Sm,Eu,Gd,Tb,Yb,Lu}, ytick={1,10,100,1000}, yticklabels={1,10,100,1000},
  ylabel={العيّنة / الكوندريت CI},
  legend style={font=\scriptsize, at={(0.98,0.98)}, anchor=north east,
  legend cell align=left, legend columns=2}]
% computed in figdata/bachelor-3/melting_models.py
\addplot[omDef, thick, mark=*, mark size=1.3pt] table[x=z, y=sp01] {figdata/out/bachelor-3-melting_models-ree.dat};
\addlegendentry{\foreignlanguage{arabic}{السبينيل، \mbox{\textdir TLT $F = 1\,\%$}}}
\addplot[red!70!black, thick, mark=square*, mark size=1.3pt] table[x=z, y=gt01] {figdata/out/bachelor-3-melting_models-ree.dat};
\addlegendentry{\foreignlanguage{arabic}{الجارنيت، \mbox{\textdir TLT $F = 1\,\%$}}}
\addplot[omDef, thick, densely dashed, mark=*, mark size=1.3pt] table[x=z, y=sp05] {figdata/out/bachelor-3-melting_models-ree.dat};
\addlegendentry{\foreignlanguage{arabic}{السبينيل، \mbox{\textdir TLT $5\,\%$}}}
\addplot[red!70!black, thick, densely dashed, mark=square*, mark size=1.3pt] table[x=z, y=gt05] {figdata/out/bachelor-3-melting_models-ree.dat};
\addlegendentry{\foreignlanguage{arabic}{الجارنيت، \mbox{\textdir TLT $5\,\%$}}}
\addplot[omDef, thick, densely dotted, mark=*, mark size=1.3pt] table[x=z, y=sp20] {figdata/out/bachelor-3-melting_models-ree.dat};
\addlegendentry{\foreignlanguage{arabic}{السبينيل، \mbox{\textdir TLT $20\,\%$}}}
\addplot[red!70!black, thick, densely dotted, mark=square*, mark size=1.3pt] table[x=z, y=gt20] {figdata/out/bachelor-3-melting_models-ree.dat};
\addlegendentry{\foreignlanguage{arabic}{الجارنيت، \mbox{\textdir TLT $20\,\%$}}}
\addplot[black, thin, mark=none] table[x=z, y=pm] {figdata/out/bachelor-3-melting_models-ree.dat};
\addlegendentry{المصدر}
\end{axis}
\end{tikzpicture}

\par\smallskip{\small أنماط العناصر الأرضية النادرة المُطبَّعة بالكوندريت لمصهورات \omterm{thm:b3:trace-element-modelling:batch}{الانصهار الدفعي}
لمصدر من الوشاح الأولي، من الليرزوليت السبينيلي (دوائر زرقاء)
ومن الليرزوليت الجارنيتي (مربعات حمراء)، عند ثلاث قيم لكسر المصهور؛ والمصدر هو
الخط الرفيع. لا تُرسم إلّا العناصر الأرضية النادرة التسعة ذات
معاملات التوزيع المجدولة.}\label{fig:b3:trace-element-modelling:ree}
\end{omfigure}
% ledger: budahn91:d-ree, ms95:pm-trace, lodders:ci-table, white13:ex73, rfr25:gkr-5gpa

يحلّ السبينيل محلّ الجارنيت بوصفه المعدن الألومينيّ الرئيس في الوشاح، على عمق يقارب 60
إلى \qty{90}{km}. فالمصهور الذي تكون عناصره الأرضية النادرة الثقيلة مسطّحة ومنخفضة، مع
انحدار حادّ من الخفيفة إلى الثقيلة، قد ترك الجارنيت في بقاياه، وبدأ
ينصهر على عمق لا يقلّ عن ذلك.
% ledger: white13:sp-gt

يمزج النمط بين مجهولين، هما $F$ والمصدر. ورسم عنصرين
أحدهما مقابل الآخر يفصل بينهما.

\begin{definition}[مخطط نسبة--نسبة]\label{def:b3:trace-element-modelling:ratios}
\emph{مخطط نسبة--نسبة}\index{مخطط نسبة--نسبة} يرسم، لمجموعة من الحمم،
نسبة عنصر شديد عدم التوافق $H$ إلى عنصر أقلّ عدم توافقًا $M$، مقابل تركيز $H$
(أو مقابل نسبة أخرى). ترسم عليه عمليات الانصهار الجزئي والتبلّر والمزج
خطوطًا ذات أشكال مختلفة
على المخطط.
\end{definition}

\begin{proposition}[خط الانصهار الدفعي]\label{prop:b3:trace-element-modelling:line}
بالنسبة إلى مصهورات دفعية من مصدر واحد، مع $D_H = 0$ و$D_M$ ثابت، تقع النقاط
$(C_H, C_H/C_M)$ على الخط المستقيم
\[
\frac{C_H}{C_M} = \frac{D_M}{C_0^M}\,C_H + \frac{C_0^H(1 - D_M)}{C_0^M}.
\]
\end{proposition}

\begin{proof}
مع $D_H = 0$، يكون $C_H = C_0^H/F$، ومنه $1/F = C_H/C_0^H$. ثمّ
$C_H/C_M = (C_0^H/F)\,(D_M + F(1 - D_M))/C_0^M = C_0^HD_M/(FC_0^M) +
C_0^H(1 - D_M)/C_0^M$، والحدّ الأوّل هو $(D_M/C_0^M)\,C_H$.
\end{proof}

يحرّك التبلّر التجزيئي للمعادن التي ترفض العنصرين عيّنةً من الحمم
أفقيًّا على هذا المخطط، لأنّ النسبة لا تتغيّر؛ ويحرّكها الانصهار
على طول خط مائل يزداد ميله مع $D_M$؛ ويرسم مزج مصهورين
قطعًا زائدًا. والمصدر الذي يحتوي على الجارنيت له $D_{\mathrm{Yb}}$ كبير، ومن ثَمّ خطّ
انصهار شديد الانحدار في مخطط La/Yb مقابل La.

\begin{method}[استنتاج من حممتين]\label{met:b3:trace-element-modelling:invert}
(1) تحقّق من أنّ الحمم متولّدة عن أصل واحد، وأنّ $H$ أشدّ عدم توافقًا
بكثير من $M$. (2) صحّح كلّ عيّنة من الحمم لأثر التبلّر، أو استعمل النسب وحدها.
(3) عبر عيّنتين أو أكثر من الحمم، طابِق الخط $C_H/C_M = aC_H + b$. (4) ثمّ
$a = D_M/C_0^M$ و$b = C_0^H(1 - D_M)/C_0^M$: بافتراض $C_0^M$، يكون $D_M =
aC_0^M$ و$C_0^H = bC_0^M/(1 - D_M)$. (5) كسر المصهور لكلّ عيّنة من الحمم
هو $F = C_0^H/C_H$.
\end{method}

\begin{omfigure}
\begin{tikzpicture}
\begin{axis}[omrecord, width=0.78\linewidth, height=0.52\linewidth,
  xmode=log, ymode=log, xmin=2, xmax=60, ymin=2, ymax=100,
  xlabel={\mbox{\textdir TLT La (\unit{ppm})}}, ylabel={\mbox{\textdir TLT La/Yb}},
  xtick={2,5,10,20,50}, xticklabels={2,5,10,20,50},
  ytick={2,5,10,20,50,100}, yticklabels={2,5,10,20,50,100},
  legend style={font=\scriptsize, at={(0.03,0.97)}, anchor=north west,
  legend cell align=left}]
% computed in figdata/bachelor-3/melting_models.py
\addplot[omDef, thick] table[x=la_sp, y=lay_sp] {figdata/out/bachelor-3-melting_models-traj.dat};
\addlegendentry{الليرزوليت السبينيلي}
\addplot[red!70!black, thick] table[x=la_gt, y=lay_gt] {figdata/out/bachelor-3-melting_models-traj.dat};
\addlegendentry{الليرزوليت الجارنيتي}
\addplot[omDef, only marks, mark=|, mark size=3pt, forget plot] table[x=la_sp, y=lay_sp] {figdata/out/bachelor-3-melting_models-ticks.dat};
\addplot[red!70!black, only marks, mark=|, mark size=3pt, forget plot] table[x=la_gt, y=lay_gt] {figdata/out/bachelor-3-melting_models-ticks.dat};
\addplot[black, only marks, mark=diamond*, mark size=2.8pt] table[x=la, y=layb] {figdata/out/bachelor-3-melting_models-lavas.dat};
\addlegendentry{\foreignlanguage{arabic}{الحمم A و\mbox{\textdir TLT B}}}
\node[font=\scriptsize, anchor=north] at (axis cs:9.6,5.2) {A};
\node[font=\scriptsize, anchor=south east] at (axis cs:15.2,36) {B};
\node[font=\scriptsize, omDef, anchor=north west] at (axis cs:36.5,17.5) {1\,\%};
\node[font=\scriptsize, omDef, anchor=north west] at (axis cs:11.3,6.2) {5\,\%};
\node[font=\scriptsize, omDef, anchor=north west] at (axis cs:3.14,2.55) {20\,\%};
\node[font=\scriptsize, red!70!black, anchor=south east] at (axis cs:29.5,66) {1\,\%};
\node[font=\scriptsize, red!70!black, anchor=south east] at (axis cs:10.5,24) {5\,\%};
\node[font=\scriptsize, red!70!black, anchor=south east] at (axis cs:3.09,7.0) {20\,\%};
\end{axis}
\end{tikzpicture}

\par\smallskip{\small مصهورات \omterm{thm:b3:trace-element-modelling:batch}{الانصهار الدفعي} لمصدر من الوشاح الأولي على مخطط La/Yb
مقابل La، للليرزوليت السبينيلي والليرزوليت الجارنيتي، مع علامات عند $F = 1$،
2، 5، 10 و\qty{20}{\percent}، والحممتين المصطنعتين لمسألة نهاية
الأسبوع. تقع كلّ عيّنة من الحمم على منحنًى واحد فقط.}\label{fig:b3:trace-element-modelling:ratios}
\end{omfigure}

على المحاور اللوغاريتمية ينحني الخطّ المستقيم الوارد في القضية، ومع
$D_{\mathrm{La}} \ne 0$ لا يكون مستقيمًا إلّا تقريبًا على المحاور الخطّية؛ أمّا منحنيات
الشكل فهي دقيقة. لا يستطيع مصدر سبينيلي أن ينتج مصهورًا تتجاوز فيه نسبة La/Yb
$(C_0^{\mathrm{La}}/D_{\mathrm{La}})/(C_0^{\mathrm{Yb}}/D_{\mathrm{Yb}})$,
، وهي نحو 40 للوشاح الأولي؛ ويمكن للمصدر الجارنيتي أن يبلغ نحو 120.

\section{التمارين}

\begin{exercise}[$\star$]\label{exo:b3:trace-element-modelling:1}
باستعمال معاملات السماريوم في المعادن، وهي 0.40 في الكلينوبيروكسين و0.11 في
الجارنيت و0.016 في الأوليفين والأورثوبيروكسين والسبينيل، احسب معامل التوزيع الكلّي للسماريوم
في الليرزوليت السبينيلي الذي يحتوي على \qty{12}{\percent} كلينوبيروكسين،
وفي الليرزوليت الجارنيتي الذي يحتوي على \qty{19.5}{\percent} كلينوبيروكسين
و\qty{16.7}{\percent} جارنيت، الواردين في الفصل.
\end{exercise}
% ledger: budahn91:d-ree

\begin{exercise}[$\star$]\label{exo:b3:trace-element-modelling:2}
عنصر له $D = 0.05$ موجود بتركيز \qty{1}{ppm} في مصدر. احسب
تركيزه في المصهورات الدفعية عند $F = 0.02$ و$0.10$.
\end{exercise}

\begin{exercise}[$\star$]\label{exo:b3:trace-element-modelling:3}
احسب أكبر إثراء $C_L/C_0$ يمكن أن يعطيه \omterm{thm:b3:trace-element-modelling:batch}{الانصهار الدفعي} لعنصر
له $D = 0.01$، ثم لعنصر له $D = 0$ عند
$F = \qty{5}{\percent}$.
\end{exercise}

\begin{exercise}[$\star$]\label{exo:b3:trace-element-modelling:4}
تحتوي حممة على \qty{30}{ppm} من اللانثانوم و\qty{2.5}{ppm} من الإيتربيوم. باستعمال قيم الكوندريت CI
للانثانوم 0.232 وللإيتربيوم \qty{0.163}{ppm}، احسب نسبتها La/Yb
ونسبتها المطبَّعة على الكوندريت $(\mathrm{La}/\mathrm{Yb})_N$.
\end{exercise}
% ledger: lodders:ci-table

\begin{exercise}[$\star\star$]\label{exo:b3:trace-element-modelling:5}
قارن، لعنصر له $D = 0.01$، البواقي بعد \qty{5}{\percent}
من \omterm{thm:b3:trace-element-modelling:batch}{الانصهار الدفعي} ومن \omterm{thm:b3:trace-element-modelling:fractional}{الانصهار التجزيئي}. ماذا يحتوي المصهور
التجزيئي اللحظي في تلك اللحظة؟
\end{exercise}

\begin{exercise}[$\star\star$]\label{exo:b3:trace-element-modelling:6}
في المثال المحلول في الفصل (النسب الشكلية 58، و27، و12، و\qty{3}{\percent}؛
نسب الانصهار 20، و25، و45، و\qty{10}{\percent} للأوليفين،
والأورثوبيروكسين والكلينوبيروكسين والسبينيل)، أمّا معاملات اللانثانوم فهي
0.0003 و0.002 و0.053 و0.01. احسب $D_0$ و$P$، ثم إثراء اللانثانوم في المصهور الدفعي
عند $F = 0.02$ و$0.10$، \omterm{def:b3:trace-element-modelling:nonmodal}{بالانصهار غير الشكلي} والانصهار الشكلي.
\end{exercise}
% ledger: white13:ex73

\begin{exercise}[$\star\star$]\label{exo:b3:trace-element-modelling:7}
عند $F = 0.10$، قارن المصهور الدفعي \omterm{def:b3:trace-element-modelling:fractional}{والمصهور التجزيئي المجمَّع} لكلٍّ من
$D = 0.1$ و$D = 0.01$. أيّهما يختلف أكثر، ولماذا؟
\end{exercise}

\begin{exercise}[$\star\star$]\label{exo:b3:trace-element-modelling:8}
عنصر له $D = 0.001$ ينصهر انصهارًا ديناميكيًا عند $\phi = 0.01$. احسب
$D'$، وإثراء أوّل مصهور مستخلَص، وإثراء المصهور المستخلَص
عند $F = 0.01$؛ وقارن ذلك \omterm{thm:b3:trace-element-modelling:fractional}{بالانصهار التجزيئي}.
\end{exercise}

\begin{exercise}[$\star\star$]\label{exo:b3:trace-element-modelling:9}
في الليرزوليت السبينيلي الوارد في الفصل، $D_{\mathrm{La}} = 0.0078$
و$D_{\mathrm{Sm}} = 0.062$. احسب العامل الذي يرفع \omterm{thm:b3:trace-element-modelling:batch}{الانصهار الدفعي} نسبة La/Sm فوق
قيمتها في المصدر عند $F = 0.01$ وعند $F = 0.10$، ثم فسّر ماذا يدلّ ذلك
عن استعمال النسب في تقدير $F$.
\end{exercise}

\begin{exercise}[$\star\star\star$]\label{exo:b3:trace-element-modelling:10}
بيّن أنّ \omterm{def:b3:trace-element-modelling:fractional}{المصهور التجزيئي المجمَّع} هو متوسط المصهورات اللحظية
من 0 إلى $F$، وأنّه يؤول إلى $C_0/F$ عندما $(1 - F)^{1/D} \to 0$.
عند $D = 0.01$، من أيّ $F$ يكون ضمن \qty{1}{\percent} من $C_0/F$؟
\end{exercise}

\begin{exercise}[$\star\star\star$]\label{exo:b3:trace-element-modelling:11}
لحمّتين متشابهتين في المنشأ، P وQ، لدينا اللانثانوم 20 و\qty{10}{ppm}، ونسبة \mbox{\textdir TLT La/Yb} تساوي 10 و6.
افترض أنّ اللانثانوم غير متوافق تمامًا، وأنّ المصدر يحتوي على
\qty{0.44}{ppm} من الإيتربيوم، وطبّق طريقة الفصل: أوجد $D_{\mathrm{Yb}}$،
واللانثانوم في المصدر، وكسر المصهور لكلّ حممة.
\end{exercise}

\begin{exercise}[$\star\star\star$]\label{exo:b3:trace-element-modelling:12}
في \omterm{def:b3:trace-element-modelling:nonmodal}{الانصهار غير الشكلي} للتمرين~6، حدّد عند أيّ كسر مصهور يُستنفد كلّ معدن،
واحسب معامل التوزيع الكلّي للبقايا من اللانثانوم عند
$F = 0.2$، واشرح لماذا لا يمكن استعمال صيغة الفصل بعد
الاستنفاد الأوّل.
\end{exercise}

\section{مسألة: حمتان من جزيرة ألدرن}

\begin{problem}[{مسألة نهاية الأسبوع --- المصدر والعمق ودرجة الانصهار
لبازلتين مصطنعتين، استنادًا إلى اللانثانوم والإيتربيوم}]\label{pb:b3:trace-element-modelling:1}
على جزيرة ألدرن، وهي جزيرة بركانية مصطنعة، تحتوي الحمم A على
\qty{9.6}{ppm} من اللانثانوم و\qty{1.7}{ppm} من الإيتربيوم، والحمم B على \qty{15.6}{ppm} من اللانثانوم و\qty{0.45}{ppm} من الإيتربيوم.
وتُعامَل الحمتان كلتاهما على أنّهما مصهورات أوّلية لمصدر من الوشاح الأولي مع
\qty{0.648}{ppm} من اللانثانوم و\qty{0.441}{ppm} من الإيتربيوم (كوندريت CI: اللانثانوم 0.232، الإيتربيوم
\qty{0.163}{ppm}). معاملات المعادن: اللانثانوم 0.058 في الكلينوبيروكسين و0.001
في الجارنيت، وفي الأوليفين والأورثوبيروكسين والسبينيل؛ والإيتربيوم 0.58 في
الكلينوبيروكسين، و4.5 في الجارنيت، و0.16 في الأوليفين والأورثوبيروكسين والسبينيل.
يتألّف الليرزوليت السبينيلي من \qty{58}{\percent} أوليفين، و\qty{27}{\percent} أورثوبيروكسين، و\qty{12}{\percent}
كلينوبيروكسين، و\qty{3}{\percent} سبينيل؛ أمّا الليرزوليت الجارنيتي فيتألّف من \qty{54.1}{\percent} أوليفين، و\qty{9.8}{\percent}
أورثوبيروكسين، و\qty{19.5}{\percent} كلينوبيروكسين، و\qty{16.7}{\percent} جارنيت. والانصهار هنا شكلي.
% ledger: ms95:pm-trace, lodders:ci-table, budahn91:d-ree, white13:ex73, rfr25:gkr-5gpa

\medskip\noindent\textbf{الجزء الأول --- المصدران.}
\begin{enumerate}
\item احسب $D_{\mathrm{La}}$ و$D_{\mathrm{Yb}}$ للليرزوليت السبينيلي.
\item احسب المعاملين نفسيهما للليرزوليت الجارنيتي. قرّب جميع القيم الأربع إلى رقمين
معنويين لبقيّة المسألة.
\item أيّ المعاملات يختلف أكثر بين المصدرين، ولماذا؟ هل
الإيتربيوم متوافق في الليرزوليت الجارنيتي؟
\item احسب $(\mathrm{La}/\mathrm{Yb})_N$ المطبَّعة على الكوندريت
للمصدر. كيف يبدو نمطه؟
\end{enumerate}

\medskip\noindent\textbf{الجزء الثاني --- \omterm{thm:b3:trace-element-modelling:batch}{الانصهار الدفعي} للحمم A.}
\begin{enumerate}[resume]
\item احسب La في المصهور الدفعي للليرزوليت السبينيلي عند $F = 1$، 5
و\qty{10}{\percent}.
\item أوجد كسر المصهور الذي يعطي الحمم A تركيزها من اللانثانوم.
\item توقّع الإيتربيوم فيها ونسبة La/Yb، وقارِن ذلك بالحمم A.
\item احسب La و Yb في الباقي، وكسر لانثانوم المصدر
المتبقّي فيه.
\item ما كسر المصهور الذي يُحصل عليه إذا عُدّ اللانثانوم
غير متوافق تمامًا؟
\end{enumerate}

\medskip\noindent\textbf{الجزء الثالث --- \omterm{thm:b3:trace-element-modelling:fractional}{الانصهار التجزيئي} للحمم A.}
\begin{enumerate}[resume]
\item احسب La في المصهور التجزيئي اللحظي عند $F = 0.1$، 1
و\qty{6}{\percent}.
\item احسب La في \omterm{def:b3:trace-element-modelling:fractional}{المصهور التجزيئي المجمَّع} عند $F = \qty{6.0}{\percent}$.
\item أوجد كسر المصهور الذي عنده يحمل المصهور المجمَّع لانثانوم الحمم A.
لماذا يقترب كثيرًا من جواب السؤال~9؟
\item قارن لانثانوم الباقي بما ورد في السؤال~8.
\item أيّ تقدير للقيمة $F$ يُعتمد؟
\end{enumerate}

\medskip\noindent\textbf{الجزء الرابع --- الحمم B والجارنيت.}
\begin{enumerate}[resume]
\item احسب La/Yb و$(\mathrm{La}/\mathrm{Yb})_N$ للحممين كليهما.
\item ما أكبر قيمة لنسبة La/Yb يمكن أن يبلغها المصهور الدفعي
من الليرزوليت السبينيلي؟
\item عند أيّ كسر مصهور يملك مصهور الليرزوليت السبينيلي نسبة La/Yb للحمم
B، وكم من اللانثانوم سيحتويه؟ استنتج.
\item أوجد كسر مصهور الحمم B انطلاقًا من لانثانومها، من أجل الليرزوليت الجارنيتي،
وتوقّع الإيتربيوم فيها.
\item ماذا يقول المصدران عن العمقين اللذين بدأ عندهما انصهار الحمم A و B،
وهل يمكن أن تنشأ الحمم A من الليرزوليت الجارنيتي؟
\item احسب ميل ومقطع خط \omterm{thm:b3:trace-element-modelling:batch}{الانصهار الدفعي} La/Yb $= a\,
\mathrm{La} + b$ للمصدر الجارنيتي باعتبار $D_{\mathrm{La}}$ معدومًا،
وتحقّق من الحمم B بمقارنتها بهذا الخط.
\item لو فقدت الحمم A \qty{10}{\percent} من كتلتها على هيئة أوليفين قبل
أن تندفع، كيف ستتغيّر نسبة La/Yb وكسر مصهورها؟
\item أجب عن السؤال الوارد في مطلع الفصل بالنسبة إلى الحمم A والحمم B.
\item حدّد درجة الانصهار للحمم A.
\end{enumerate}
\end{problem}
